% =====================================================================
%  Chapter 16 — The Cosmological Tensions
% =====================================================================

\chapter{The Cosmological Tensions}
\label{ch:tensions}

The strongest reason to think that cosmology is about to produce a discovery
is not that it has found one. It is that the standard model of cosmology is
under measurable strain in three places at once, and the strains are of a
kind that a dark matter component could relieve.

\section{The Hubble tension}The CMB-inferred value of $H_0$ is
$67.36 \pm 0.54\ \kms$ \idx{Planck Collaboration}. The local measurements
using Cepheid variables, Type Ia supernovae and masers give values closer to
73 to 74, and the tension between the two is larger than the combination of
the quoted uncertainties suggests.

The tension has been under pressure from both sides for years. The
measurement of the CMB itself improved; the local ladder improved; the
distance to the Large Magellanic Cloud, the key calibration rung, has been
revised; and the 2024 JWST measurements of Cepheids in the SN Ia hosts
confirmed the higher local values rather than lowering them. The
discrepancy is now large enough that it is generally described as requiring
new physics rather than as an unresolved systematic, though the review
literature of 2025 and 2026 continues to catalogue the remaining
possibilities \idx{Hubble tension review}.

The relevant point for this book is narrow and specific: the tension is in
the \emph{expansion rate today}, and a dark matter component that is not
exactly pressureless --- a slightly interacting or a decaying species ---
modifies the expansion history in the right direction. It is one of several
reasons the field is looking hard at the dark matter sector's less standard
corners.

\section{\texorpdfstring{$S_8$}{S8} tension}A second tension concerns the amplitude of density
fluctuations. The CMB measures $\sigma_8 = 0.811$, and the combination with
the baryon density gives $S_8 \equiv \sigma_8 (\Omega_m/0.3)^{0.6}$. Weak
lensing and galaxy clustering measurements have tended to prefer a lower
value, by a few per cent.

The significance of the $S_8$ tension is lower than the Hubble tension and its
status is more contested: the two are difficult to compare directly because
they probe different quantities and different redshift ranges, and the
redshift evolution of the tension has been argued over extensively. But the
direction is the same as the Hubble tension, and the combination of the two
is the reason that modified-gravity and interacting-dark-matter models remain
part of the live discussion despite the failure of the modified-gravity
programme at the cluster scale.

\section{The DESI result}The new element of 2025 was the DESI second data
release, which combined baryon acoustic oscillation measurements from more
than 14 million galaxies with CMB and supernova data and found that the
combination preferred an evolving dark energy equation of state over the
standard $\Lambda$CDM at between 2.8 and 4.2 sigma, depending on the data
combination used \idx{DESI Collaboration}.

The result was the most discussed cosmological measurement of the decade,
and it was discussed for the right reasons: the significance is high enough
to matter and low enough that it can vanish, the conclusion depends on the
supernova sample, and the data release is public.

The subsequent re-analysis by the DES-Dovekie group reduced the maximum
significance to about 3.2 sigma by re-examining the modelling choices in the
supernova data, which is exactly the kind of check that should follow such a
result \idx{DES-Dovekie reanalysis}. A 3.2 sigma preference is not a
discovery. It is a statement that the question is now open.

\section{The supernova bias}One of the mechanisms proposed to reconcile the
DESI result with standard cosmology is a systematic bias in the supernova
sample, arising because the progenitor population evolves: younger
progenitors are brighter, and a population whose progenitors age over the
redshift interval biases the luminosity distance --- and hence the inferred
dark energy equation of state --- at the few per cent level. A 2025
analysis in the Monthly Notices quantified the effect at a level comparable
to the DESI signal itself \idx{SN progenitor age}.

This is a striking example of a systematic that has been argued about for
years, at the level of a few per cent, in a direction that matters, and which
would now have to be measured with a precision comparable to the effect it is
being invoked to explain.

\section{Why this matters for dark matter}Three reasons, in increasing order of
importance.

The first is that a cosmological discovery of evolving dark energy would be
the largest change in the field since the discovery of the accelerating
expansion, and it would happen in the same decade in which the dark matter
particle is expected. The two would be hard to disentangle, and the
possibility that they are the same phenomenon --- a dark sector that
modifies both the expansion and the clustering --- is a live and testable
idea.

The second is that the tensions are the strongest available probe of dark
matter interactions at scales no laboratory can reach. If dark matter has a
small but non-zero coupling to the Standard Model, it can relax the Hubble
tension while remaining invisible in the direct-detection experiments.

The third is strategic. Chapter~\ref{ch:roadmap} argues that the field should
be preparing for the possibility that the cosmological route produces the
first evidence, and that this possibility is currently under-resourced
relative to the laboratory programme.

\begin{keytakeaways}
\item The Hubble tension, the $S_8$ tension and the DESI dark energy
result are three distinct strains in the standard model of cosmology, and
they are of a kind that a dark sector component could relieve.
\item The DESI result stands at 2.8 to 4.2 sigma depending on the data
combination, and a re-analysis reduced the maximum to about 3.2 sigma. It is
an open question, not a discovery.
\item A progenitor-age bias in the supernova sample is of the same size as
the DESI signal, which is precisely why it must now be measured.
\item The cosmological route may deliver the first evidence for new physics
before the laboratory does, and the field is not resourced for that.
\end{keytakeaways}
